\section{System Architecture}
\label{sec:architecture}

\subsection{High-Level Overview}

The system consists of four main components:

\begin{enumerate}
  \item \textbf{Training Pipeline}: Loads the MovieLens 100K dataset, trains the SVD model, and persists it.
  \item \textbf{Model Artifact}: The trained, serialized model file (\code{models/svd\_model.pkl}).
  \item \textbf{Serving API}: A FastAPI application that loads the model and handles HTTP requests.
  \item \textbf{Client Interaction}: External clients communicate with the API via REST over HTTP(S).
\end{enumerate}

The workflow is:

\begin{verbatim}
  [MovieLens 100K Data]
           |
           v
  [Training Script (train_model.py)]
           |
           v
  [SVD Model Artifact (svd_model.pkl)]
           |
           v
  [API Startup - Load Model]
           |
           v
  [FastAPI Application Running]
           |
           v
  [Client Requests] --> [Health Check]
                     --> [Single Prediction]
                     --> [Batch Predictions]
\end{verbatim}

\subsection{Training Pipeline}

\subsubsection{Data Loading}

The \code{scripts/train\_model.py} script begins by loading the MovieLens 100K dataset:

\lstinputlisting[language=Python, linerange={25,35}]{../../scripts/train_model.py}

The Surprise library handles downloading and caching the dataset automatically. The \code{prompt=False} 
parameter ensures non-interactive execution (important for CI/CD environments).

\subsubsection{Cross-Validation}

Before training on the full dataset, the script performs 5-fold cross-validation to assess model generalization:

\begin{verbatim}
  - Split data into 5 folds
  - For each fold: train on 4 folds, evaluate on 1 fold
  - Compute mean and variance of RMSE and MAE across folds
  - Use results to inform hyperparameter selection
\end{verbatim}

Cross-validation provides confidence that the model generalizes beyond the training set and is not 
severely overfitting.

\subsubsection{Full Dataset Training}

After cross-validation, the model is trained on the \textit{entire} dataset:

\lstinputlisting[language=Python, linerange={45,55}]{../../scripts/train_model.py}

Training on the full dataset maximizes the model's exposure to available data for production deployment. 
This is a standard practice in ML: use cross-validation for evaluation during development, then train on 
all data for the final artifact.

\subsubsection{Model Serialization}

The trained model is serialized and persisted:

\lstinputlisting[language=Python, linerange={60,68}]{../../scripts/train_model.py}

Python's \code{pickle} module is used for simplicity. For production systems handling sensitive data, 
more secure serialization formats might be considered.

\subsection{API Architecture}

\subsubsection{Startup and Model Loading}

When the FastAPI application starts, it loads the model in the startup event handler:

\begin{verbatim}
  @app.on_event("startup")
  async def startup_event():
      global model
      try:
          model = MovieRatingModel()
          logger.info("Model loaded successfully at startup")
      except Exception as e:
          logger.error(f"Failed to load model: {e}")
\end{verbatim}

This design ensures the model is loaded once at startup, then reused for all subsequent requests. 
If model loading fails, the application continues but runs unhealthy (reported by the health check endpoint).

\subsubsection{Model Wrapper Class}

The \code{app/model.py} module encapsulates model logic in a \code{MovieRatingModel} class:

\begin{itemize}
  \item \code{\_\_init\_\_(model\_path)}: Load the pickled model from disk.
  \item \code{predict(user\_id, movie\_id)}: Single prediction with clamping to $[1.0, 5.0]$.
  \item \code{predict\_batch(pairs)}: Batch predictions (wrapper around single predictions).
  \item \code{is\_loaded()}: Check if model is successfully loaded.
\end{itemize}

This encapsulation separates model logic from API concerns, improving testability and maintainability.

\subsubsection{Request/Response Pipeline}

The API flow for a prediction request is:

\begin{enumerate}
  \item Client sends JSON POST request with \code{user\_id} and \code{movie\_id}.
  \item FastAPI validates the request using Pydantic schemas (Section~\ref{sec:api}).
  \item If validation fails, return 422 (Unprocessable Entity).
  \item The endpoint handler checks if the model is loaded.
  \item If model is not loaded, return 503 (Service Unavailable).
  \item Model \code{predict()} method is called.
  \item Response object is constructed and serialized to JSON.
  \item Client receives 200 OK with prediction.
\end{enumerate}

\subsection{Containerization}

The application is containerized using Docker to ensure reproducibility across environments:

\begin{enumerate}
  \item \textbf{Base Image}: \code{python:3.10-slim} provides a minimal Python environment.
  \item \textbf{Dependencies}: System packages (curl for healthchecks) and Python packages are installed.
  \item \textbf{Code and Models}: Application code and model artifact are copied into the container.
  \item \textbf{Health Checks}: Docker HEALTHCHECK ensures the container remains healthy.
  \item \textbf{Startup Command}: \code{uvicorn} is configured to run the FastAPI application.
\end{enumerate}

Docker Compose orchestrates the deployment:

\begin{verbatim}
  version: '3.8'
  services:
    api:
      build: .
      ports:
        - "8000:8000"
      volumes:
        - ./models:/app/models
      environment:
        - MODEL_PATH=/app/models/svd_model.pkl
      restart: unless-stopped
\end{verbatim}

The \code{volumes} directive mounts the local \code{./models} directory into the container, enabling 
model artifact updates without container rebuilds.

\subsection{Deployment Scenarios}

The system supports three deployment modes:

\textbf{Local Development:}
\begin{verbatim}
  pip install -r requirements.txt
  python scripts/train_model.py
  uvicorn app.main:app --reload --host 0.0.0.0 --port 8000
\end{verbatim}

\textbf{Local Production (No Docker):}
\begin{verbatim}
  python scripts/train_model.py
  uvicorn app.main:app --host 0.0.0.0 --port 8000
\end{verbatim}

\textbf{Containerized (Docker Compose):}
\begin{verbatim}
  docker compose build
  docker compose up -d
  curl http://localhost:8000/health
\end{verbatim}

Each mode uses the same application code, ensuring consistency.
